\section{Preliminares}

Podíamos definir \(\| T(x) \| \le \| T \|_{\infty} \cdot \| x \| \quad \forall x \in \mathbb{R}^n\).

\begin{theorem}[Regla de la Cadena]
    Sea \(A \subset \mathbb{R}^n\) abierto, \(F: A \subset \mathbb{R}^n \to \mathbb{R}^m\), y \(G: B \subset \mathbb{R}^m \to \mathbb{R}^k\) con \(F(A) \subset B\). Si \(F\) es diferenciable en \(p \in A\) y \(G\) es diferenciable en \(F(p) \in B\), entonces \(G \circ F\) es diferenciable en \(p\) y además:
    \[ D(G \circ F)_p = DG_{F(p)} \cdot DF_p \]
\end{theorem}

\subsection{Teoremas de Análisis I}

\begin{theorem}[Teorema de Fermat]
    Sea \(f: [a,b] \to \mathbb{R}\). Supongamos que \(c \in (a,b)\) es extremo local de \(f\). Si \(f\) es derivable en \(c\), entonces \(f'(c) = 0\).
\end{theorem}

\begin{theorem}[Teorema de Rolle]
    Si \(f\) es continua en \([a,b]\), derivable en \((a,b)\) y \(f(a) = f(b)\), entonces \(\exists c \in (a,b)\) tal que \(f'(c) = 0\).
\end{theorem}

\begin{theorem}[Teorema de Lagrange]
    Sea \(f\) continua en \([a,b]\) y derivable en \((a,b)\). Entonces \(\exists c \in (a,b)\) tal que:
    \[ f'(c) = \frac{f(b) - f(a)}{b - a} \]
\end{theorem}

\subsection{Extensión a \(\mathbb{R}^n\)}

\begin{theorem}[Lagrange en \(\mathbb{R}^n\)]
    Sea \(f: B_r(p) \subset \mathbb{R}^n \to \mathbb{R}\) diferenciable. Para todo \(a, b \in B_r(p)\), en el segmento que une \(a\) con \(b\), existe \(c\) tal que:
    \[ f(b) - f(a) = \nabla f(c) \cdot (b - a) \]
\end{theorem}

\begin{theorem}[Teorema de Fermat en \(\mathbb{R}^n\)]
    Sea \(f: A \subset \mathbb{R}^n \to \mathbb{R}\) con \(A\) abierto. Si \(p \in A\) es extremo local de \(f\), entonces \(\nabla f(p) = 0\).
\end{theorem}

\subsection{Propiedades de Transformaciones Lineales}
\begin{enumerate}
    \item Toda transformación lineal \(T\) es diferenciable y \(DT_p = T\) para todo \(p \in \mathbb{R}^n\).
    \item Si \(T: \mathbb{R}^m \to \mathbb{R}^k\) es lineal y \(F \in C^1\), entonces \(T \circ F\) es diferenciable y \(D(T \circ F)_p = T \cdot DF_p\).
    \item \(T\) es inyectiva si y solo si \(N(T) = \{0\}\).
    \item Si \(T: \mathbb{R}^n \to \mathbb{R}^n\), se tiene que \(T\) es inversible si y solo si es inyectiva (por teorema de la dimensión).
\end{enumerate}

\begin{lemma}
    Sea \(G: A \subset \mathbb{R}^n \to \mathbb{R}^m\) abierto. Si \(G\) es \(C^1\) en \(A\), dados \(p, q \in A\), vale que:
    \[ \| DG_p - DG_q \|_{\infty} \le n \max_{i,j} \left| \frac{\partial g_i}{\partial x_j}(p) - \frac{\partial g_i}{\partial x_j}(q) \right| \]
\end{lemma}

\begin{lemma}
    Sea \(G: B_r(p) \subset \mathbb{R}^n \to \mathbb{R}^m\) diferenciable y supongamos que \(\| DG_x \|_{\infty} \le M\) para todo \(x \in B_r(p)\). Entonces:
    \[ \| G(x) - G(y) \| \le M \| x - y \| \quad \forall x,y \in B_r(p) \]
\end{lemma}

\section{Teorema de la Función Inversa}

\begin{theorem}[Teorema de la Función Inversa]
    Sea \(F: A \subset \mathbb{R}^n \to \mathbb{R}^n\) con \(A\) abierto. Si \(F\) es \(C^1\) en \(A\) y \(DF_p: \mathbb{R}^n \to \mathbb{R}^n\) es una transformación lineal inversible para un \(p \in A\), entonces:
    \begin{enumerate}
        \item \(\exists r > 0\) tal que \(F: B_r(p) \to F(B_r(p))\) es biyectiva y la función inversa es continua en su dominio.
        \item Existen abiertos \(U \subset B_r(p)\) y \(W \subset F(B_r(p))\) tales que \(F^{-1}: W \to U\) es diferenciable.
        \item La diferencial de la función inversa es \(D(F^{-1})_{F(q)} = (DF_q)^{-1}\).
        \item Si además \(F \in C^k(U)\), entonces \(F^{-1} \in C^k(W)\).
    \end{enumerate}
\end{theorem}
\begin{proof}
    1) Supongamos primero que \(DF_p = I\). Usando la continuidad de la diferencial, existe \(r_0 > 0\) tal que \(B_{r_0}(p) \subset A\). Tomamos \(r < r_0\) tal que \(\|DG_x - DF_p\|_\infty \le \frac{1}{2}\) para \(\|x-p\| < r\).
    Definimos \(G(x) = F(x) - x\), entonces \(DG_x = DF_x - I\), y \(\|DG_x\|_\infty \le \frac{1}{2}\). Por el lema anterior, \(\|G(x) - G(y)\| \le \frac{1}{2}\|x-y\|\).
    Desarrollando \(\|x - y\|\) tenemos:
    \[ \|x - y\| = \|x - y - (F(x) - F(y)) + (F(x) - F(y))\| \le \|G(x) - G(y)\| + \|F(x) - F(y)\| \]
    Sustituyendo la cota de \(G\):
    \[ \|x - y\| \le \frac{1}{2}\|x-y\| + \|F(x) - F(y)\| \implies \frac{1}{2}\|x-y\| \le \|F(x) - F(y)\| \]
    Esto prueba que si \(x \neq y\), \(F(x) \neq F(y)\), por ende \(F\) es inyectiva en \(B_r(p)\) y su inversa es continua (es bi-lipschitz).

    2) Demostremos que cubre un abierto. Sea \(S_r(p) = \{x \in \mathbb{R}^n : \|x-p\| = r\} \subset A\). \(S_r(p)\) es compacto. Para un \(y\) en \(W = B_{d/2}(F(p))\) (donde \(d\) es la distancia de \(F(p)\) a \(F(S_r(p))\)), consideramos \(h(x) = \|y - F(x)\|^2\).
    Al ser continua en el compacto \(\overline{B_r}(p)\), alcanza su mínimo en un punto \(p_y\). Evaluando en los bordes se verifica que \(p_y\) debe estar en el interior \(B_r(p)\). Al ser mínimo local diferenciable, evaluando en cualquier dirección \(v\):
    \[ \langle v, DF_{p_y}^t (y - F(p_y)) \rangle = 0 \]
    Como el operador diferencial es inversible por estar en la bola, \(DF_{p_y}^t\) también lo es, deduciéndose que \(y - F(p_y) = 0 \implies F(p_y) = y\).

    3) La diferenciabilidad se deduce mediante una manipulación del límite \(z \to F(q)\) apoyándonos en que \(F\) es diferenciable y que su inversa local acaba de ser probada como continua.

    Finalmente, si \(DF_p \neq I\), tomamos \(\tilde{F} = (DF_p)^{-1} \circ F\) que verifica \(D\tilde{F}_p = I\) y hereda las propiedades demostradas, extendiendo el resultado al caso general.
\end{proof}

\begin{lemma}
    Sean \(A, B: \mathbb{R}^n \to \mathbb{R}^n\) transformaciones lineales.
    \begin{enumerate}
        \item Si \(A\) es inversible y \(\| B - A \|_{\infty} < \frac{1}{\| A^{-1} \|_{\infty}}\), entonces \(B\) es inversible.
        \item El conjunto de operadores lineales inversibles en \(\mathbb{R}^n\) es abierto, y la función \(A \mapsto A^{-1}\) es continua.
    \end{enumerate}
\end{lemma}
\begin{proof}
    1) Llamemos \(\frac{1}{\alpha} = \|A^{-1}\|_\infty\) y \(\beta = \|B-A\|_\infty < \alpha\). Sea \(x \in \mathbb{R}^n\).
    Por un lado, \(\|x\| = \|A^{-1}Ax\| \le \frac{1}{\alpha}\|Ax\| \implies \alpha\|x\| \le \|Ax\|\).
    Además, \(\|Ax\| = \|(A-B+B)x\| \le \|(A-B)x\| + \|Bx\| \le \beta\|x\| + \|Bx\|\).
    Uniendo ambas desigualdades:
    \[ \alpha\|x\| \le \beta\|x\| + \|Bx\| \implies (\alpha-\beta)\|x\| \le \|Bx\| \]
    Como \(\beta < \alpha\), tenemos \(\alpha - \beta > 0\). Si \(x \neq 0\), la desigualdad implica \(\|Bx\| > 0 \implies Bx \neq 0\). Por lo tanto, \(N(B) = \{0\}\), lo que significa que \(B\) es inyectiva. En dimensión finita, esto equivale a ser inversible.

    2) Tomando \(x = B^{-1}y\) en la desigualdad anterior obtenemos \((\alpha - \beta)\|B^{-1}y\| \le \|y\|\), lo que nos da que \(\|B^{-1}\|_\infty \le \frac{1}{\alpha-\beta}\). Para ver la continuidad, usamos la identidad \(B^{-1} - A^{-1} = B^{-1}(A-B)A^{-1}\).
    Tomando normas:
    \[ \|B^{-1} - A^{-1}\|_\infty \le \|B^{-1}\|_\infty \|A-B\|_\infty \|A^{-1}\|_\infty \]
    Esto tiende a 0 a medida que \(B \to A\), probando que la aplicación toma la inversa de manera continua.
\end{proof}